\section{Discussion and limitations}\label{sec:discussion}

Making every kernel a nonnegative combination of phase-type densities turns Hawkes
estimation into a certified convex problem and turns responses to shocks into matrix
exponentials; adding a neural history encoder and a defective hyper-Erlang renewal channel to
the same state space gives a model whose likelihood stays exact and which is competitive with,
and on order books and Retweet better than, the strongest neural temporal point processes.
Three lessons stand out. First, absorption speed is not a property of the news kernel: echoes
lengthen it by one to two orders of magnitude, and the correction is available in closed form.
Second, falsification matters: two of our specifications produced plausible-looking news
effects that the placebo test exposed as artefacts. Third, benchmark comparisons must check
that every model scores the same observed data---a clamp inside one implementation was enough
to reverse a ranking.

\paragraph{Limitations.} Expectations for payrolls, claims, PPI and retail sales are
statistical rather than survey-based, and the FX feed has a gap in the first half of 2023. The
order-book evaluation uses the public LOBSTER samples (one trading day per stock). Neural
baselines on order books were trained with reduced batch sizes on an 8\,GB GPU, and on the
public benchmarks we compare against published numbers, validated by one re-run. On Retweet the
EPT-X budget was 25 epochs with the validation likelihood still rising, so the reported margin
is if anything conservative. EPT-X does not reach the best published likelihood on Taxi, where
the gap lies in mark prediction, or on Amazon, where it lies in the timing density; and the
MSX forecasts of post-release activity are beaten by regressions fitted to that target.
